\documentclass[11pt]{ctexart}
\usepackage{amsmath,amssymb,longtable}
\begin{document}
\section*{P415: 当前完整数学模型（2026-10-20）}
\section*{修订依据与边界}补齐源汇场站时间及整个FIFO旅行函数；固定发车0，有限状态标签遍历所有可行请求顺序，整数路线选择与完整修复模型等价，没有启发式标签截断。

\section*{完整时变取送路线模型}
请求集合$R$，取货$i$与送货$i+|R|$；源汇$0,2|R|+1$；最多$K$台车。
每车路线从源到汇，不含自环、不重复客户，出发时刻0；每请求可拒绝，
接受则同车先取后送且只一次。服务开始$s_i\in[e_i,l_i]$，
离开$d_i=s_i+\sigma_i$，到达$s_j\ge d_i+\tau_{ij}(d_i)$。
所选弧包括源汇弧全部满足这一传递。载重从0开始，取货增$q_i$、送货减$q_i$，全程在$[0,Q]$内并空载返回。
时变旅行时长唯一满足
\[
\int_t^{t+\tau_{ij}(t)}v(u)\,du=dist_{ij}.
\]
$v(u)>0$为实例7个完整时区的速度，跨时区连续积分，最后速度继续覆盖规划期；因此FIFO成立。
\[
\max \sum_{i\text{被服务}}profit_i-c_t\sum_{k\text{使用}}return_k-Z\,\#\{\text{使用车辆}\}.
\]
\section*{等价有限路线选择形式}
对每个可行服务子集$A$，枚举全部取送顺序取得最早返回$T_A$。
标签为(已取集合,已送集合,末节点)，只删除同标签更晚离开者：
因允许等待、FIFO以及非负时间成本，早标签可模拟所有晚标签后续，故这是精确支配。
\[
z_A\in\{0,1\},\quad\sum_Az_A\le K,\quad\sum_{A\ni i}z_A\le1,\quad
\max\sum_A\left(\sum_{i\in A}profit_i-c_tT_A-Z\right)z_A.
\]
无标签数/路线数启发式截断。全部4组数据乘5种速度情景独立求解，目标不跨实例相加。

\end{document}
